\section{Equilibrium Analysis}
\label{sec:analysis}

We characterise the equilibria of the Privacy Externality Game in five
moves. \Cref{lem:externality} shows the mechanism induces a monotone
externality on the observable $A_i$, and
\cref{cor:externality-small-p,thm:externality-full} promote its sign from
an empirical observation on the trade-off function to an analytical
theorem on the strict interior of the feasible region. Stated in slack
form (\cref{prop:slack-monotonicity}), the same fact says the DP-worst-case
gap shrinks whenever another player tightens their budget.
\Cref{prop:no-symmetric-ne} then rules the diagonal out as a Nash
equilibrium in the fully symmetric two-player game. On the small-$p$
sublattice \cref{thm:continuous-ne-smallp} delivers pure-strategy NE
existence on the continuous action space via the supermodular route of
\citet{milgrom1990rationalizability}; the discretised game grounds our
numerical search. \Cref{prop:non-potential} rules out any exact,
weighted, or ordinal potential representation, explaining why the search
proceeds by enumeration rather than by minimising a scalar functional.
\Cref{prop:V-invariance} isolates the load-bearing role of $V$ (only the
sign of $V'$ enters), and \cref{cor:coalition-attack} identifies the
budget-manipulation attack of \citet{kaiser2026your} as the rational
best response of any coordinated subset of data holders.

\subsection{Monotonicity of the externality}

\begin{lemma}[Privacy externality on the feasible region]
\label{lem:externality}
Consider the sampling-based iDP mechanism of \cref{def:sampling-idp} on a
feasible profile $\boldsymbol{\eps}$ at which constraints (C1) bind for
every group and (C2) binds at $E_b$. Let $\sigma$ and $p_g$ denote the
mechanism's solution. Then for every pair $i \neq j$,
\begin{equation}
\frac{\partial \sigma}{\partial \eps_j} \;<\; 0
\qquad\textrm{and}\qquad
\frac{\partial p_i}{\partial \eps_j} \;<\; 0.
\label{eq:externality-lemma}
\end{equation}
That is, when player $j$ loosens its budget, the mechanism uses less noise
($\sigma$ falls) and reduces group $i$'s sampling rate ($p_i$ falls).
Analytically (\cref{cor:externality-small-p,thm:externality-full}), the
realised MIA advantage $A_i$ also falls on the strict interior of the
feasible region: $\partial A_i / \partial \eps_j < 0$.
\end{lemma}

\begin{proof}
Use the R\'enyi-DP form of (C1) of \citet{mironov2017renyi}: at order
$\alpha$, the RDP of an $I$-fold $(p_g,\sigma)$-subsampled Gaussian is at
most $\rho_g \approx 2I\alpha\, p_g^2/\sigma^2$ (tight up to higher-order
terms for $p_g \ll 1$). Setting this equal to the target RDP
$\rho_g(\eps_g, \delta)$ (which is monotone-\emph{increasing} in $\eps_g$
because $\rho = \eps - \log(1/\delta)/(\alpha-1)$ for the binding RDP-to-DP
conversion at any chosen order $\alpha > 1$) gives the per-group sampling
rate
\[
p_g(\sigma, \eps_g) \;=\; \frac{\sigma}{\sqrt{2I\alpha}}\,
\sqrt{\rho_g(\eps_g, \delta)}.
\]
Substituting into (C2) yields
\begin{equation}
\sigma \cdot \sum_{g}\, |G_g|\, \sqrt{\rho_g(\eps_g,\delta)} \;=\; E_b
\sqrt{2I\alpha}.
\label{eq:sigma-implicit}
\end{equation}
Implicit differentiation in $\eps_j$ gives
$\partial \sigma / \partial \eps_j = -\sigma\, |G_j|\, (\partial \sqrt{\rho_j}
/ \partial \eps_j) / [\sum_g |G_g|\sqrt{\rho_g}] < 0$, since
$\rho_j$ is increasing in $\eps_j$. From the closed-form expression for
$p_i$ at fixed $\eps_i$, $\partial p_i / \partial \eps_j = (\partial \sigma /
\partial \eps_j) \cdot \sqrt{\rho_i(\eps_i)}/\sqrt{2I\alpha} < 0$.

The observable sign $\partial A_i / \partial \eps_j < 0$ chains through
$\partial A_i / \partial p_i > 0$ and $\partial A_i / \partial \sigma < 0$
of the $I$-fold $(p_i,\sigma)$-subsampled Gaussian, which pull in
opposite directions (more sampling exposes more; more noise hides more).
The dominance of the $p_i$-decrease effect is established analytically:
\cref{cor:externality-small-p} closes the sign in the small-$p$ regime by
a leading-order expansion of the subsampled trade-off function, and
\cref{thm:externality-full} extends the sign to the strict interior of
the feasible region via a Chernoff two-sided envelope on the composed
privacy-loss distribution.
\end{proof}

\Cref{lem:externality} formalises the empirical observation of
\citet{kaiser2026your}: when any player tightens their iDP budget, the
mechanism's noise multiplier $\sigma$ shifts to keep (C2) binding, and every
other player's realised vulnerability $A_i$ rises. The externality is
\emph{monotone} in every other player's budget.

\begin{corollary}[Externality on $A_i$ in the small-$p$ regime]
\label{cor:externality-small-p}
Under the hypotheses of \cref{lem:externality}, assume additionally that
$p_g \leq p_\star$ for every group $g$, with $p_\star$ small enough that
the $O(p^2)$ remainder in the amplification expansion of
\citet{balle2020privacy} is dominated by the leading term uniformly in
$\alpha \in [0,1]$. Then
\begin{align}
A_i(p_i,\sigma)
    &= p_i\, \Adv_G(\sigma) \;+\; O(p_i^2),
        \label{eq:cor-Ai-leading}\\
\frac{\partial^2 A_i}{\partial p_i\, \partial \sigma}
    &= \Adv_G'(\sigma) \;+\; O(p_i) \;<\; 0,
        \label{eq:cor-mixed}
\end{align}
where $\Adv_G(\sigma) = 2\Phi(1/(2\sigma))-1$ is the MIA advantage of the
unsubsampled sensitivity-1 Gaussian. Consequently,
$\partial A_i / \partial \eps_j < 0$ on the small-$p$ face of the
feasible region.
\end{corollary}

\begin{proof}[Proof sketch]
The Poisson-subsampled trade-off function admits the mixture representation
$T_{p,\sigma}(\alpha) = (1-p)(1-\alpha) + p\, T_\sigma(\alpha) + O(p^2)$
(\cite{balle2020privacy}, Poisson-subsampling amplification), so its
sup-gap $A(p,\sigma) = \sup_\alpha[1-\alpha-T_{p,\sigma}(\alpha)]$ equals
$p\,\Adv_G(\sigma) + O(p^2)$; \eqref{eq:cor-Ai-leading} and
\eqref{eq:cor-mixed} follow by term-by-term differentiation. For the
externality sign, substitute the closed-form $p_i = \sigma\sqrt{\rho_i}/
\sqrt{2I\alpha}$ from the proof of \cref{lem:externality} into the chain
rule for $\partial A_i/\partial\eps_j$ and factor
$(\partial\sigma/\partial\eps_j)<0$ out:
\begin{equation}
\frac{\partial A_i}{\partial \eps_j}
   \;=\; \frac{\partial \sigma}{\partial \eps_j}
        \cdot p_i \cdot
        \Bigl[\, \tfrac{\Adv_G(\sigma)}{\sigma}
              + \Adv_G'(\sigma)\, \Bigr] \;+\; O(p_i^2).
\label{eq:proof-bracket}
\end{equation}
The bracket is the $\sigma$-derivative of $\sigma\,\Adv_G(\sigma)$
divided by $\sigma$; setting $u := 1/(2\sigma)$ and
$f(u) := \Phi(u) - \Phi(-u) - 2u\phi(u)$ gives $f(0)=0$ and
$f'(u) = 2u^2\phi(u) > 0$, so the bracket is strictly positive. Together
with $(\partial\sigma/\partial\eps_j)<0$, the leading $O(p_i)$ term of
$\partial A_i/\partial\eps_j$ is strictly negative.
\end{proof}

\begin{theorem}[Full-region monotonicity of $A_i$]
\label{thm:externality-full}
Under the hypotheses of \cref{lem:externality}, assume the strict-interior
condition $p_g(\boldsymbol{\eps}) \in (0,1)$ for every group $g$. Then for
every pair $i \neq j$, $\partial A_i / \partial \eps_j < 0$ on the
strict interior of the feasible region.
\end{theorem}

\begin{proof}[Proof sketch]
The composed Poisson-subsampled Gaussian trade-off function
$T_{p,\sigma}$ is sandwiched between two Gaussian trade-off functions
\[
T_{\mathcal G(\mu_+(p,\sigma))}(\alpha) \;\leq\; T_{p,\sigma}(\alpha)
   \;\leq\; T_{\mathcal G(\mu_-(p,\sigma))}(\alpha),
\qquad \alpha \in [0,1],
\]
with envelope scales
$\mu_+(p,\sigma) = 2p\sqrt{I}/\sigma$
(RDP-to-zCDP-to-GDP conversion; \cite{mironov2017renyi,bun2016concentrated})
and
$\mu_-(p,\sigma) = p\sqrt{I}\sqrt{e^{1/\sigma^2}-1}$
(GDP-CLT lower envelope; \cite{dong2022gaussian}, Cor.~3.3).
Both scales are strictly increasing in $p$ and strictly decreasing in
$\sigma$, and $\Adv(\mathcal G(\mu)) = 2\Phi(\mu/2)-1$ is strictly
increasing in $\mu$. Chaining through
$\partial\sigma/\partial\eps_j < 0$ and $\partial p_i/\partial\eps_j<0$
of \cref{lem:externality} on the strict interior, both envelopes carry
$\partial A_i / \partial \eps_j < 0$; the true $A_i$ inherits the sign
from the Balle-Barthe-Gaboardi privacy profile
(\cite{balle2020privacy}, Thm.~9), which is strictly increasing in
$p$ and strictly decreasing in $\sigma$ at every fixed $\eps > 0$.
\end{proof}

In the decomposition of \cref{sec:decomposition}, the lemma equivalently
states that \emph{the slack $\bar c(\eps_i) - A_i(\boldsymbol{\eps})$
shrinks in every other player's $\eps_j$}. In words, the
$(\eps_i,\delta)$-DP worst-case ceiling $\bar c(\eps_i)$ is unchanged
(it depends only on $\eps_i$), but the realised advantage $A_i$ climbs
toward that ceiling as others tighten their budgets. We record this
slack-form as a proposition because the $(\eps_i, \delta, \Delta)$-iDP
contract of \citet{kaiser2026your, kaissis2024beyond} caps precisely this
quantity, and every downstream reference to the externality's sign can
cite a single labelled statement.

\begin{proposition}[Slack monotonicity]
\label{prop:slack-monotonicity}
Fix a profile $\boldsymbol{\eps}$ in the feasible region of
\cref{ass:participation} at which the hypotheses of \cref{lem:externality}
hold. Then for every pair $i \neq j$,
\[
\frac{\partial \mathrm{slack}_i}{\partial \eps_j}(\boldsymbol{\eps})
\;=\; -\,\frac{\partial A_i}{\partial \eps_j}(\boldsymbol{\eps})
\;>\; 0
\]
on the strict interior. Loosening any other player's budget strictly
enlarges $i$'s slack (pushes $A_i$ away from its $(\eps_i,\delta)$-DP
ceiling); tightening any other player's budget strictly shrinks it.
\end{proposition}

\begin{proof}
Differentiate the identity $A_i(\boldsymbol{\eps}) = \bar c(\eps_i)
- \mathrm{slack}_i(\boldsymbol{\eps})$. Because $\bar c(\eps_i)$ depends
only on $\eps_i$, $\partial \bar c / \partial \eps_j = 0$ for $j\neq i$
and the equality $\partial \mathrm{slack}_i / \partial \eps_j =
-\,\partial A_i / \partial \eps_j$ is immediate. The strict inequality
is \cref{thm:externality-full} on the strict interior of the feasible
region, upgraded from the small-$p$ closure of
\cref{cor:externality-small-p}.
\end{proof}

\paragraph{Minority and majority intuition.}
The externality has an asymmetric \emph{group-size} structure. Because (C2)
enforces a fixed expected batch size, a player with small $|G_i|$ but
small $\eps_i$ contributes little to the batch and forces the
larger-$|G_j|$ players' sampling rates upward. The minority therefore
\emph{exports} its privacy demand onto the majority. This is the formal
content of the empirical observation in \cite{kaiser2026your}
that ``the proportion of low-budget players matters as much as their
$\eps$''.

\subsection{No symmetric equilibrium in the symmetric 2-player PEG}
\label{sec:no-symmetric-ne}

\Cref{lem:externality} has an immediate structural consequence for the
fully symmetric two-player game: the diagonal cannot be a pure Nash
equilibrium. This is, together with \cref{lem:externality}, the
analytical content of our headline empirical finding (\cref{fig:C1}) that
all pure NEs of the symmetric 2-player PEG sit off the diagonal.

\begin{proposition}[No symmetric pure NE in the symmetric 2-player PEG]
\label{prop:no-symmetric-ne}
Consider the two-player PEG of \cref{def:peg} with $n = 2$, identical
privacy weights $w_1 = w_2 = w > 0$, and identical group sizes
$|G_1| = |G_2|$. Assume the standing regularity of \cref{eq:V} on $V$ and
that the vulnerability $A_i$, viewed as a function of $(p_i,\sigma)$, is
$C^2$ with $\partial A / \partial p > 0$ strictly on the interior of the
feasible region. Let $(\eps^\star, \eps^\star)$ be an interior diagonal
profile at which both mechanism constraints (C1) and (C2) hold with strict
feasibility $p_i(\eps^\star, \eps^\star) > 0$. Then $(\eps^\star,
\eps^\star)$ is \emph{not} a pure Nash equilibrium.
\end{proposition}

\begin{proof}
\emph{Setup.} Change variables to $s_i := \sqrt{\rho(\eps_i, \delta)}$
where $\rho(\cdot,\delta)$ is the RDP-order-$\alpha$ bound from the proof
of \cref{lem:externality}. Then $\eps_i \mapsto s_i$ is a $C^1$
diffeomorphism on the interior; substituting the closed form of
$p_g(\sigma,\eps_g)$ into (C2) yields, with $|G_1| = |G_2| = m$,
\begin{equation}
\sigma(s_1, s_2) \;=\; \frac{C}{s_1 + s_2},
\qquad
p_i(s_1, s_2) \;=\; \frac{C}{\sqrt{2 I \alpha}}\cdot\frac{s_i}{s_1 + s_2},
\label{eq:agg-form}
\end{equation}
with $C := E_b \sqrt{2 I \alpha}/m > 0$. Write $S := s_1 + s_2$; utility
becomes $u_i(s_1, s_2) = V(\sigma(S)) - w\, A(p_i, \sigma(S))$, with the
same $A(\cdot,\cdot)$ for both players by the symmetric group sizes.

\emph{Aggregator structure at the diagonal.} Fix $s_1 = s_2 = s > 0$, so
$S = 2s$, $\sigma = C/(2s)$, and $p_1 = p_2 = C/(2\sqrt{2 I\alpha}) =: p$.
Write $a := \partial A / \partial p |_{(p,\sigma)}$ and
$b := \partial A / \partial \sigma |_{(p,\sigma)}$; by hypothesis $a > 0$.
From \eqref{eq:agg-form},
\[
\frac{\partial \sigma}{\partial s_i} \;=\; -\frac{\sigma}{2 s},
\qquad
\frac{\partial p_1}{\partial s_1} \;=\; \frac{p}{2s},
\qquad
\frac{\partial p_1}{\partial s_2} \;=\; -\frac{p}{2s}.
\]
The sign asymmetry $\partial p_1 / \partial s_1 > 0 > \partial p_1 /
\partial s_2$ is the mechanism's aggregator content: player~$i$'s
sampling rate is monotone in her share $s_i/S$, so raising her own
$s_i$ has the opposite effect on $p_1$ as raising her opponent's $s_j$.

\emph{Diagonal FOC collapses.} Player~1's first-order condition in $s_1$
is
\begin{equation}
\frac{\partial u_1}{\partial s_1}
\;=\; \bigl[V'(\sigma) - w b\bigr] \Bigl(-\tfrac{\sigma}{2 s}\Bigr)
    \;-\; w a \,\tfrac{p}{2 s}
\;=\; 0.
\label{eq:foc-diagonal}
\end{equation}
By the symmetric relabelling $1\leftrightarrow 2$, $\partial u_2 /
\partial s_2$ at $(s,s)$ is the identical expression: the two FOCs
coincide, and any diagonal candidate is characterised by the one scalar
equation \eqref{eq:foc-diagonal}.

\emph{Sign of the cross-partial.} Player~1's marginal utility with respect
to \emph{player~2's} action at the diagonal is
\[
\frac{\partial u_1}{\partial s_2}
\;=\; \frac{\partial u_1}{\partial s_1}
   \;+\; w a\, \Bigl(\frac{\partial p_1}{\partial s_1} - \frac{\partial
   p_1}{\partial s_2}\Bigr)
\;=\; 0 \;+\; w a \cdot \frac{p}{s}
\;>\; 0,
\]
using \eqref{eq:foc-diagonal} in the first line and the aggregator
identity in the second. The strict inequality relies only on $w,a,p,s>0$.

\emph{Second-order test in own action.} Because $\sigma$ depends on
$(s_1,s_2)$ only through the aggregate $S = s_1 + s_2$, the shared-$V$
and $\sigma$-derivative contributions to $\partial^2 u_1/\partial s_1^2$
and $\partial^2 u_1 / (\partial s_1 \partial s_2)$ are identical and
cancel in their difference. Only the $p_1$-derivative terms survive; a
direct computation from \eqref{eq:agg-form} (writing $a_{pp} :=
\partial^2 A / \partial p^2$ at $(p,\sigma)$) gives, at the diagonal,
\begin{equation}
\frac{\partial^2 u_1}{\partial s_1^2}
\;-\;
\frac{\partial^2 u_1}{\partial s_1 \partial s_2}
\;=\; \frac{w\, p}{2\, s^2}\,\bigl(a \;-\; a_{pp}\, p\bigr).
\label{eq:soc-difference}
\end{equation}
The mixed partial $\partial^2 u_1 / (\partial s_1 \partial s_2)$ is the
derivative along $s_1$ of the strictly positive quantity $\partial u_1
/ \partial s_2 = w a p / s$ evaluated at the FOC; the aggregator
representation makes it non-negative at $(s^\star, s^\star)$ (raising
$s_1$ towards $s_2$'s value leaves $\partial u_1 / \partial s_2$
weakly increasing on the diagonal ray). Under the sign condition
\begin{equation}
a \;-\; a_{pp}\, p \;>\; 0
\qquad\text{at }(p,\sigma) \;=\; \bigl(C/(2\sqrt{2I\alpha}),\, C/(2 s^\star)\bigr),
\label{eq:sign-condition}
\end{equation}
\eqref{eq:soc-difference} is strictly positive, so $\partial^2 u_1 /
\partial s_1^2 > 0$ strictly at the diagonal FOC: the profile $(s^\star,
s^\star)$ is a strict interior \emph{minimum} of $u_1(\cdot, s^\star)$
in $s_1$, not a maximum.

Condition \eqref{eq:sign-condition} asserts that the vulnerability
$A(p,\sigma)$ is more responsive to $p$ than a log-linear (constant
elasticity) function would predict at the diagonal operating point. It
holds throughout the mechanism-solver caches used in
\cref{sec:experiments} (verified numerically on every diagonal profile
of the C1 and C2 grids and symbolically for the leading-order
$A(p,\sigma) = p\,\Adv_G(\sigma) + O(p^2)$ expansion of
\cref{cor:externality-small-p}, where $a - a_{pp}\,p = \Adv_G(\sigma) +
O(p) > 0$).

\emph{Conclusion.} At the diagonal FOC candidate $(s^\star, s^\star)$,
$\partial u_1/\partial s_1 = 0$ and $\partial^2 u_1/\partial s_1^2 > 0$
strictly, so
\[
u_1(s^\star + t, s^\star) - u_1(s^\star, s^\star)
\;=\; \tfrac{1}{2}\,\frac{\partial^2 u_1}{\partial s_1^2}\bigg|_{(s^\star, s^\star)}\, t^2
\;+\; O(t^3)
\;>\; 0
\]
for all sufficiently small $t \neq 0$: player~1 has a strictly
profitable unilateral deviation in either direction. Non-FOC diagonal
profiles fail the first-order test directly. By the $C^1$
diffeomorphism $\eps_i \mapsto s_i$, no interior diagonal profile
$(\eps^\star, \eps^\star)$ is a pure Nash equilibrium.
\end{proof}

\Cref{prop:no-symmetric-ne} is a purely negative statement: the diagonal
is ruled out. It does not claim that a pure NE exists at any asymmetric
profile; positive existence on the continuous action space is the content
of \cref{thm:continuous-ne-smallp} below. The proposition also does not
rely on the specific closed form $V(\sigma) = 1/(1+\sigma^2)$: the only
property of $V$ used is $V'(\sigma) < 0$ strictly, so the analogous
statement holds for every $V$ satisfying \cref{ass:V-admissible}.

\subsection{Existence and structure of pure Nash equilibria}
\label{sec:existence}

\paragraph{Scaffolding: existence on the discretised game.}
The discretised PEG on the log-spaced action grid $\mathcal{A}_K
\subset [\eps^{\min},\eps^{\max}]$ (\cref{sec:experiments}) with
utilities~\eqref{eq:utility} is a finite normal-form game restricted to
the (closed) feasible region of \cref{ass:participation}, so a
mixed-strategy NE exists by~\cite{nash1950equilibrium}. This grounds the
numerical search: the exhaustive-enumeration procedure of
\cref{sec:experiments} has at least one fixed point to find, and
restricting attention to $\mathcal{A}_K$ is warranted. Pure NE on
$\mathcal{A}_K$ need not exist and depend on $(w_i, |G_i|, V, K)$; we
enumerate them by brute force and verify stability under grid refinement.

Existence on the continuous action space is not automatic---the utility
\eqref{eq:utility} is not generally concave in $\eps_i$ because $A_i$ is
the tight trade-off function of an $I$-fold subsampled Gaussian, whose
dependence on $\eps_i$ through $(\sigma, p_i)$ has no closed form, so the
standard Debreu-Fan-Glicksberg quasi-concavity route does not apply.
The order-theoretic route of \citet{topkis1979equilibrium,
milgrom1990rationalizability} does, on the small-$p$ sublattice:

\begin{theorem}[Pure-strategy NE existence on the small-$p$ sublattice]
\label{thm:continuous-ne-smallp}
Let $\mathcal F_\star \subseteq \mathcal F$ be the closure of the set of
feasible profiles at which $p_g(\boldsymbol{\eps}) \leq p_\star$ and
$\sigma(\boldsymbol{\eps}) \geq 1/\sqrt 3$ for every group $g$, with
$p_\star$ chosen as in \cref{cor:externality-small-p}. Under the
reversed-$\eps$ product order, $\mathcal F_\star$ is a compact sublattice
of $\Sigma$, and the $n$-player PEG restricted to $\mathcal F_\star$
admits a non-empty complete lattice of pure-strategy Nash equilibria; in
particular there exist a largest and a smallest NE
$\boldsymbol{\eps}^{\mathrm{sup}}, \boldsymbol{\eps}^{\mathrm{inf}}
\in \mathcal F_\star$.
\end{theorem}

\begin{proof}[Proof sketch]
Expand $\partial^2 u_i / (\partial \eps_i \partial \eps_j)$ by chain
rule; the shared-$V$ contribution
$V''(\sigma)\,\sigma_i\sigma_j + V'(\sigma)\,\sigma_{ij}$ carries a
non-negative sign once $\sigma \geq 1/\sqrt 3$, while the $-w_i A_i$
contribution's dominant $O(1)$ terms in the small-$p$ regime are
$-w_i A_{i,p\sigma}\,(p_{i,i}\sigma_j + p_{i,j}\sigma_i)$ with
$A_{i,p\sigma} = \Adv_G'(\sigma) + O(p_i) < 0$
(\cref{cor:externality-small-p}). The bracket has the sign of
$\Adv_G(\sigma)/\sigma + \Adv_G'(\sigma)$, negative on the interior
(\eqref{eq:proof-bracket}), so the full term is negative and dominates
subleading $O(p_i)$ contributions. Therefore
$\partial^2 u_i / (\partial \eps_i \partial \eps_j) \leq 0$ on
$\mathcal F_\star$, i.e.\ $u_i$ has increasing differences under the
reversed-$\eps$ order. The mechanism-solver monotonicity of \cref{lem:externality}
in each coordinate gives that componentwise meet and join in the
reversed-$\eps$ order preserve feasibility, so $\mathcal F_\star$ is a
compact sublattice. Milgrom-Roberts Theorem~5
\cite{milgrom1990rationalizability} then delivers a non-empty complete
lattice of pure NEs with a largest and a smallest element.
\end{proof}

The condition $\sigma \geq 1/\sqrt 3$ is a mild restriction (the K=200
mechanism-grid cache confirms $\sigma \in [0.5, 5]$ throughout the
paper's operating region), and the small-$p$ sublattice covers the
DP-SGD-motivated sampling rates $p \sim 10^{-4}$ to $10^{-2}$ used in
every scenario reported in \cref{sec:experiments}. Existence on the
full closed feasible region (including the $p_g \to 1$ saturation corner
and the feasibility edge) is not proven analytically; we rely on the
exhaustive grid enumeration of \cref{sec:experiments} to certify the
remaining region.

\paragraph{What the externality lemma implies for NE structure.}
\Cref{lem:externality} says lowering another player's budget weakly
raises one's own realised vulnerability. If a player's utility places
non-trivial weight on $A_i$ ($w_i > 0$), she has an incentive to
\emph{shift the externality onto others} by lowering her own budget---an
incentive that is intrinsic to the mechanism. The opposing incentive comes
from $V(\sigma)$: lowering one's own $\eps$ pushes $\sigma$ upward,
lowering the model utility for everyone (herself included). The
equilibrium balances these two incentives.

\Cref{prop:no-symmetric-ne} pins the qualitative structure of that
balance in the symmetric two-player game: no interior profile
$(\eps,\eps)$ is a pure NE. Our numerical evaluation (\cref{sec:experiments})
confirms that the pure NEs found by exhaustive search all sit off the
diagonal, and are located at profiles where one player carries a higher
$A$ in exchange for a higher $V$-contribution.

\paragraph{Locus of NEs.}
The set of pure NEs of the 2-player game can be characterised implicitly
through the first-order condition of each player's best response. Writing
out the partial derivative of~\eqref{eq:utility} and using the
parametrisation of the mechanism in the proof of \cref{lem:externality},
\[
\frac{\partial u_i}{\partial \eps_i} \;=\;
V'(\sigma)\, \frac{\partial \sigma}{\partial \eps_i}
\;-\; w_i\, \frac{\partial A_i}{\partial \eps_i}
\;=\; 0,
\]
i.e.\ at the BR
\begin{equation}
\frac{1}{w_i}
\;=\; \frac{ (\partial A_i / \partial \eps_i) }
          { V'(\sigma)\,(\partial \sigma / \partial \eps_i) }
\;\;\;\Big|_{\boldsymbol{\eps}}.
\label{eq:foc}
\end{equation}
For each fixed $w_i$, \cref{eq:foc} defines an implicit curve
$\mathrm{BR}_i$ in the $(\eps_1, \eps_2)$-plane; the pure NEs are the
intersections of $\mathrm{BR}_1$ and $\mathrm{BR}_2$. The two quantities
on the right-hand side of~\eqref{eq:foc} are not available in closed form
(the dependence of $A_i$ on $(\sigma, p_i)$ has no closed form), but each
is computable numerically by a single PLD-accounting evaluation. We use
this in \cref{sec:experiments} to trace both BR loci on a fine grid and
identify their intersections.

\paragraph{Non-potential structure.}
The Pejo et al.\ Price-of-Privacy game~\cite{pejo2019together} was a
potential game because its privacy cost was separable across players,
reducing equilibrium computation to minimising a single scalar potential.
The PEG's cost, coupled through the mechanism's shared $(\sigma, p_i)$,
breaks that reduction.

\begin{proposition}[Non-potential game]
\label{prop:non-potential}
Under the utility form \eqref{eq:utility} and the mechanism of
\cref{def:sampling-idp}, the PEG is neither an exact, weighted, nor
ordinal potential game in the sense of \citet{monderer1996potential}:
\begin{enumerate}
  \item[\rm(i)] there exist no positive weights $\{\eta_i\}_{i=1}^n$
        and no twice-differentiable $\Phi:\prod_i\Sigma_i \to \R$
        satisfying $u_i(\eps_i,\boldsymbol{\eps}_{-i}) -
        u_i(\eps_i',\boldsymbol{\eps}_{-i}) = \eta_i\bigl(
        \Phi(\eps_i,\boldsymbol{\eps}_{-i}) -
        \Phi(\eps_i',\boldsymbol{\eps}_{-i})\bigr)$
        on the interior of the feasible region;
  \item[\rm(ii)] and there exists a length-four cycle of strict
        unilateral improvements on the discretised feasible grid,
        precluding any ordinal potential (\cite{monderer1996potential},
        Lemma~2.5).
\end{enumerate}
\end{proposition}

\begin{proof}[Proof sketch]
\emph{(i) Cardinal potentials.} Implicit differentiation of
\eqref{eq:sigma-implicit} gives
$\sigma_{ij} := \partial^2\sigma/(\partial\eps_i\partial\eps_j)
= 2\sigma\,|G_i||G_j|(r')^2/(rN)^2$
at the symmetric-budget reference, where $r := \sqrt{\rho}$ and
$N := \sum_g|G_g|$; the shared-$V$ contribution
$V''(\sigma)\sigma_i\sigma_j + V'(\sigma)\sigma_{ij}$ is therefore
symmetric under $i\leftrightarrow j$ and cancels on both sides of the
Monderer-Shapley cross-partial condition
$\eta_j\, \partial^2 u_i / \partial\eps_i\partial\eps_j
= \eta_i\, \partial^2 u_j / \partial\eps_i\partial\eps_j$. What remains
is the $A_i$-contribution, which, after expanding
$p_{i,ij} = (\sigma_{ij} r_i + \sigma_j r_i')/\sqrt{2 I \alpha}$ and
substituting, reduces the residual to
$\eta_j w_i |G_j|(2|G_i| - N) = \eta_i w_j |G_i|(2|G_j| - N)$ at every
symmetric profile. For heterogeneous groups $|G_i|\neq|G_j|$ this
equation has no solution with positive $(\eta_i, \eta_j)$; for the fully
symmetric case the residual carries the $V$-contribution
$(\eta_j - \eta_i)\,d^2 V$, which is nonconstant along the feasible
interval and cannot be zeroed uniformly by any positive
$(\eta_1,\eta_2)$ with $\eta_1 \neq \eta_2$.

\emph{(ii) Ordinal potentials.} By \citet{monderer1996potential},
Lemma~2.5, a game admits an ordinal potential iff it has no strict
improvement cycle. On the C1 discretisation ($w_1=w_2=1$, $|G_1|=|G_2|=5000$,
$K=22$ log-spaced budget grid on $[0.5,64]$), the profile sequence
$(\eps_1,\eps_2) = (1.587,5.040) \to (2.000,5.040) \to (2.000,4.000)
\to (1.587,4.000) \to (1.587,5.040)$ has each step strictly improving
for the deviating player (unilateral-utility gaps $7.3\!\times\!10^{-7}$,
$1.5\!\times\!10^{-3}$, $3.0\!\times\!10^{-3}$, $2.8\!\times\!10^{-3}$),
so no ordinal potential exists. Together, (i) and (ii) rule out the
hierarchy of exact, weighted, and ordinal potentials.
\end{proof}

\Cref{prop:non-potential} explains why our equilibrium analysis proceeds
by discretised NE search (\cref{sec:experiments}) rather than by
minimising a closed-form potential: neither Monderer-Shapley
best-response dynamics nor the finite-improvement property is available.
The existence route we do use---order-theoretic supermodularity on the
small-$p$ sublattice (\cref{thm:continuous-ne-smallp})---trades concavity
for monotonicity, which is why the sign analysis of
$\partial A_i / \partial \eps_j$
(\cref{cor:externality-small-p,thm:externality-full})
is the load-bearing input.

\subsection{Coalition Stackelberg recovers the attack}

\Cref{lem:externality} applied to the coalition variant of
\cref{def:coalition} pins the malice-maximising coalition's Stackelberg
action down to a single point: the participation-preserving minimum
budget. We record this as a corollary.

\begin{corollary}[Coalition best response recovers the attack]
\label{cor:coalition-attack}
Fix a target $t \in N$, a coalition $C \subseteq N \setminus \{t\}$ of
size $k$, and any profile $\boldsymbol{\eps}_{-C} \in \prod_{j \notin C}
\Sigma_j$ for the non-coalition players. Let $\lambda \to \infty$ in
\cref{def:coalition}, so the coalition maximises $A_t$ alone. Then the
coalition's Stackelberg best response is the coalition-symmetric profile
$\eps_i = \eps^{\star}_k(\boldsymbol{\eps}_{-C})$ for every $i \in C$,
where $\eps^{\star}_k(\boldsymbol{\eps}_{-C})$ is the infimum of the
coalition-symmetric budgets at which the induced profile is feasible in
the sense of \cref{ass:participation}. Furthermore, holding
$\boldsymbol{\eps}_{-C}$ fixed, the induced target advantage
$A_t(\eps^{\star}_k, \boldsymbol{\eps}_{-C})$ is non-decreasing in $k$.
\end{corollary}

\begin{proof}[Proof sketch]
By \cref{lem:externality} (strict form on the interior via
\cref{thm:externality-full}), $\partial A_t / \partial \eps_i < 0$ for
every $i \in C$ on the feasible region. Setting every $\eps_i$ as low as
the feasibility region admits therefore maximises $A_t$. The
participation-preservation qualifier is necessary because at extremely
low $\eps_i$ the mechanism may set $p_i = 0$ (the colluder effectively
drops out, removing its constraint from (C2) and so its externality on
the target); the coalition's optimum is thus the infimum of feasible
$\eps_i$. Monotonicity in $k$ is the same lemma applied to adding a
member to $C$: doing so strictly enlarges the set of budgets that pull
$\sigma$ downward, raising $A_t$.
\end{proof}

\Cref{cor:coalition-attack} identifies the budget-manipulation attack
of \citet{kaiser2026your} as the coalition's equilibrium best response:
its contribution is the game-theoretic mapping, not an empirical
strengthening of that work. Whether the coalition is adversarial is
irrelevant, since any group of aware players whose joint utility weights
their target's harm heavily picks the same action vector as a rational
best response.

\subsection{Qualitative invariance of the model-utility function}
\label{sec:V-invariance}

The equilibrium results above are developed with the specific proxy
$V(\sigma) = 1/(1+\sigma^2)$ (\cref{eq:V}), motivated by the convex-DP-ERM
excess-risk bound of \citet{bassily2014private}. A natural question is
whether they hinge on that closed form. They do not: only the sign of
$V'$, not its shape or magnitude, enters the load-bearing arguments. We
isolate the load-bearing class.

\begin{assumption}[Admissible model-utility function]
\label{ass:V-admissible}
Let $\mathcal S := \sigma(\mathcal F) \subset \R_{>0}$ be the range of
noise multipliers taken by the mechanism on the feasible region. A
function $V \colon \mathcal S \to \R$ is \emph{admissible} iff
\begin{enumerate}
\item[\rm(A1)] $V$ is continuous on $\mathcal S$;
\item[\rm(A2)] $V$ is strictly decreasing on $\mathcal S$;
\item[\rm(A3)] $V$ is bounded on $\mathcal S$.
\end{enumerate}
\end{assumption}

\begin{proposition}[Qualitative invariance in $V$]
\label{prop:V-invariance}
Under \cref{ass:V-admissible}, for every admissible $V$:
\begin{enumerate}
\item[\rm(i)] the externality signs of \cref{lem:externality},
      \cref{cor:externality-small-p}, and \cref{thm:externality-full}
      hold on the feasible region;
\item[\rm(ii)] the coalition-Stackelberg conclusion of
      \cref{cor:coalition-attack} holds; and
\item[\rm(iii)] at any interior individual best response, the FOC
      \eqref{eq:foc} reduces to the $V$-free sign identity
      $\mathrm{sgn}(\partial A_i/\partial \eps_i)
       = \mathrm{sgn}(V'(\sigma)\,\partial \sigma/\partial \eps_i)$.
\end{enumerate}
\end{proposition}

\begin{proof}
The mechanism's constraints (C1)+(C2) and the closed-form $p_g(\sigma,\eps_g)$
do not read $V$: every step of the proof of \cref{lem:externality}, and of
\cref{cor:externality-small-p} and \cref{thm:externality-full} which
chain through it, is a manipulation of the mechanism's constraints and
the per-group RDP. Clause (i) follows. Clause (ii) uses only clause (i)
applied to the target's coordinate; $V$ does not enter the
$\lambda\to\infty$ coalition objective. For clause (iii), differentiate
$u_i = V(\sigma) - w_i A_i$ in $\eps_i$: by (A2), $V'(\sigma) \leq 0$
almost everywhere on $\mathcal S$ (Lebesgue's theorem on monotone
functions); by (i), $\partial\sigma/\partial\eps_i < 0$ on the feasible
interior. Setting $\partial u_i / \partial \eps_i = 0$ gives the FOC
\eqref{eq:foc}, whose sign structure depends on $V$ only through
$\mathrm{sgn}\,V'$. Real-valuedness of $u_i$ on the finite action grid
follows from (A1) and (A3), so the discretised-game scaffolding of
\cref{sec:existence} applies to every admissible $V$.
\end{proof}

\Cref{prop:V-invariance} isolates the load-bearing role of $V$: only the
sign of $V'$, not its magnitude or shape, enters the externality lemma
and the coalition-attack recovery. The experimental choice $V(\sigma) =
1/(1+\sigma^2)$ is therefore a directional proxy for a whole class of
admissible model-utility functions, and the equilibrium phenomena
reported in \cref{sec:experiments} are properties of the class, not of
the representative.

\subsection{An aside on awareness}

\Cref{lem:externality} requires that players compute (or approximate) the
realised $A_i$. A natural question is: what about a player whose only
information is the formal $(\eps_i, \delta)$-iDP contract? Such a player
treats $\eps_i$ as a direct privacy proxy and plays the budget that
maximises $V(\sigma) - g(\eps_i)$ for some monotone $g$. Their best response
does not depend on $\eps_{-i}$ at all (the externality is invisible), so
their action is mechanically the smallest $\eps_i$ that does not break
their own $V$. This is the \emph{rational} action of a naive player, and it
is precisely the configuration in which the attack of
\citet{kaiser2026your} succeeds against them: by \cref{lem:externality},
naive players' choices feed the externality on everyone else without their
realising. The contribution of this paper is to show that even
\emph{aware} players---those who can compute $A_i$ explicitly---arrive at
asymmetric equilibria that concentrate vulnerability on a subset of players.
The externality cannot be undone by individual rationality alone.
